% =====================================================================
%  1bat-ud2-14.tex  —  HORA 14: Repàs global i síntesi (estacions)
%  (sense preàmbul: s'inclou des de main.tex)
% =====================================================================
\horatitol{Sessió 14 — Repàs global i síntesi (estacions)}%
{Repassem tota la unitat rotant per quatre estacions: llenguatge, valor numèric, operacions i arrels.}

\subsection{Idea clau}

A la sessió 13 vam tancar el bloc d'arrels comprovant-les gràficament amb
GeoGebra. Avui no apareix cap contingut nou: la sessió és un \textbf{repàs
actiu} de tot el recorregut de la unitat, organitzat en quatre estacions
que roten cada $\approx$\,10 minuts. Cada estació posa a prova un dels
quatre blocs treballats.

\begin{idea}
\textbf{El mapa de la unitat, en quatre passos.}
\begin{itemize}[leftmargin=1.4em, itemsep=2pt]
  \item \textbf{Llenguatge} (sessions 1–3, criteri C1): traduir una
        situació real a una expressió algebraica o un polinomi.
  \item \textbf{Valor numèric} (sessions 4–5, criteri C2): substituir
        valors concrets i calcular, respectant la jerarquia
        d'operacions.
  \item \textbf{Operacions} (sessions 6–9, criteri C3): sumar, restar i
        multiplicar polinomis per combinar o ampliar pressupostos.
  \item \textbf{Arrels} (sessions 10–13, criteri C4): igualar a zero i
        trobar el llindar on una prestació s'exhaureix.
\end{itemize}
Del llenguatge a la fórmula, de la fórmula al valor numèric i a les
operacions, i finalment a les arrels i el seu significat: aquest és el
fil que connecta tota la unitat.
\end{idea}

\subsection*{Estació A — Llenguatge algebraic}

\begin{exercicis}[Tasca de l'estació A]
\begin{exlist}
  \item Un servei de suport escolar cobra $8\eur$ per sessió més una
        matrícula fixa de $25\eur$. Tradueix a una expressió algebraica
        el cost total per a $x$ sessions.

  \item L'expressió $12x^2-5x+40$ modelitza el cost d'un programa de
        lleure per a $x$ participants. Indica'n el grau, el coeficient
        del terme de primer grau i el terme independent.
\end{exlist}
\end{exercicis}

\subsection*{Estació B — Valor numèric}

\begin{exercicis}[Tasca de l'estació B]
\begin{exlist}[start=3]
  \item \textbf{(Compte amb la jerarquia)} El copagament d'un servei
        d'atenció domiciliària es calcula amb
        $C(x)=\num{0,05}x^2-10x+500$, on $x$ és la renda mensual en
        euros. Calcula $C(60)$ i interpreta el resultat.

  \item Un programa de suport calcula un índex de vulnerabilitat amb
        $V(x,n)=300n-x$, on $x$ són els ingressos mensuals de la
        família (en euros) i $n$ el nombre de membres. Calcula
        $V(900,4)$.
\end{exlist}
\end{exercicis}

\subsection*{Estació C — Operacions amb polinomis}

\begin{exercicis}[Tasca de l'estació C]
\begin{exlist}[start=5]
  \item Dos casals comparteixen una activitat conjunta. El cost d'un és
        $A(x)=2x^2+5x+20$ i el de l'altre $B(x)=x^2-3x+15$, on $x$ és el
        nombre d'infants inscrits. Calcula el cost conjunt $A(x)+B(x)$.

  \item El cost net d'un torn de «Tallers d'estiu» del casal de joves
        és $M(x)=50x^2+300x$ (com a la sessió 7), on $x$ és el nombre de
        joves atesos. Calcula i simplifica el cost de tres torns
        idèntics, $3\per M(x)$.
\end{exlist}
\end{exercicis}

\subsection*{Estació D — Arrels}

\begin{exercicis}[Tasca de l'estació D]
\begin{exlist}[start=7]
  \item Una prestació es modelitza amb $P(x)=-15x+450$, on $x$ són els
        ingressos en centenars d'euros. Troba el llindar d'exclusió
        (l'arrel) i interpreta'l.

  \item \textbf{(Interpretació)} Un servei calcula el seu cost net amb
        $D(x)=x^2-10x+21$. Troba les arrels amb la fórmula general i
        digues quants llindars té aquest servei.
\end{exlist}
\end{exercicis}

\bigskip
\noindent Un cop completada la rotació, posem en comú alguna resposta de
cada estació i construïm entre tots un petit mapa de la unitat com el de
la idea clau d'avui, anotant quins blocs generen més dubtes: aquests
seran precisament els que reforçarem a la sessió 15, el taller de
preparació de la prova.

% ---------------------------------------------------------------------
\fullsolucions{Sessió 14}
\begin{solucions}
\begin{sollist}
  \item $8x+25$ (en euros, per a $x$ sessions).
  \item Grau $2$; coeficient del terme de primer grau: $-5$; terme
        independent: $40$.
  \item $C(60)=\num{0,05}\per3600-10\per60+500=180-600+500=80$: el
        copagament és $80\eur$.
  \item $V(900,4)=300\per4-900=1200-900=300$.
  \item $A(x)+B(x)=3x^2+2x+35$.
  \item $3\per M(x)=150x^2+900x$.
  \item $-15x+450=0 \Rightarrow x=30$ (centenars d'euros, és a dir
        $3.000\eur$): a partir d'aquest nivell d'ingressos ja no hi ha
        dret a la prestació.
  \item $\Delta=(-10)^2-4\per1\per21=100-84=16$, $\sqrt{\Delta}=4$;
        $x=\dfrac{10\pm4}{2} \Rightarrow x_1=7$, $x_2=3$. El servei té
        \textbf{dos} llindars (el cost net s'anul·la amb $3$ i amb $7$
        unitats).
\end{sollist}
\end{solucions}
